% Decision item 1: targeted replacements for AI_Agent_For_Eye_Disease-12.pdf.
% This is an insertion document, not the full manuscript source.
% Scope: FairVision statistical reporting only. Other revision items remain open.
% Numbers: outputs/audits/fairvision_statistics_revision_20261005/.
% All predictions are historical. No API calls, new test cases, or new agents.

% ABSTRACT, page 1: replace the FairVision performance sentence only.
On the complete 250-case FairVision glaucoma and diabetic-retinopathy subsets,
RetinAgent had higher macro-F1 point estimates than RETFound (0.771 versus
0.748 and 0.860 versus 0.840), but paired confidence intervals for the
differences included zero. The AMD analysis was restricted to 210 valid
responses among 250 selected cases and did not establish full-cohort performance.

% INTRODUCTION, page 3: replace the repeated FairVision superiority sentence.
We compared RetinAgent with foundation models and standalone language models
using saved FairVision predictions. On the complete glaucoma and DR subsets,
macro-F1 differences versus RETFound were 0.0228 (95\% CI, $-0.0274$ to
0.0743) and 0.0201 ($-0.0081$ to 0.0486), respectively. These estimates do
not establish superiority or equivalence. AMD comparisons were restricted
to valid responses and are reported separately from complete-cohort results.

% RESULTS 2.3, pages 7--8: replace the heading and its three paragraphs.
\subsection{Diagnostic performance on the FairVision subsets}
RetinAgent produced valid saved predictions for all 250 glaucoma cases
(124 positive, 126 negative) and all 250 DR cases (125 positive, 125 negative).
For glaucoma, macro-F1 was 0.7712 (95\% CI, 0.7160--0.8229), compared with
0.7484 for RETFound. RetinAgent identified 89 of 124 positive cases and
104 of 126 negative cases, whereas RETFound identified 79 and 109,
respectively. The paired macro-F1 difference was 0.0228 (95\% CI,
$-0.0274$ to 0.0743). RetinAgent corrected 23 RETFound errors but introduced
18 errors on cases RETFound classified correctly; the exact two-sided
McNemar $p$ value was 0.5327 (Holm-adjusted $p=1.000$).

For DR, RetinAgent's macro-F1 was 0.8597 (95\% CI, 0.8157--0.9000),
compared with 0.8396 for RETFound. RetinAgent had 102 true positives,
113 true negatives, 12 false positives and 23 false negatives. The paired
macro-F1 difference was 0.0201 (95\% CI, $-0.0081$ to 0.0486), with nine
corrections and four regressions relative to RETFound (exact McNemar
$p=0.2668$; Holm-adjusted $p=1.000$). Thus, neither complete-cohort
comparison established an advantage over RETFound at the available precision.
Marginal intervals for all table metrics are supplied in the accompanying
interval tables; multiplicity-adjusted macro-F1 difference intervals are
retained in the statistical supplement.

For AMD, 40 of the 250 saved predictions were invalid, leaving 210 valid
responses (114 positive and 96 negative references). Among these responses,
there were 101 true positives, 85 true negatives, 11 false positives and
13 false negatives. Macro-F1 was 0.8850 (95\% CI, 0.8416--0.9278),
sensitivity was 0.8860 and specificity was 0.8854. The previously reported
0.8858 was support-weighted F1, not macro-F1. On the same 210 cases,
the macro-F1 difference versus the available VisionFM OCT predictions was
$-0.0092$ (95\% CI, $-0.0579$ to 0.0429); versus URFound OCT it was
0.0279 ($-0.0245$ to 0.0805). Neither comparison supported superiority.
These are response-conditioned comparisons, not estimates for all 250 cases.
VisionFM OCT had the highest macro-F1 among the available saved foundation
predictions on the full locked AMD subset, but those saved VisionFM and
URFound rows differ from the earlier manuscript values. Their source-version
discrepancy remains unresolved; the present table identifies the files used
and does not certify their identity with the earlier manuscript runs.

Worst-group fixed-label macro-F1 was 0.5729 for RetinAgent versus 0.5341
for RETFound in glaucoma, and 0.7795 versus 0.7582 in DR. These are descriptive
point estimates, not evidence that subgroup disparities were reduced.
RetinAgent's AMD worst-group macro-F1 was 0.4375 among valid responses.
The contributing Black subgroup contained only two positive and 16 negative
references, with neither positive case detected. The earlier worst-group
column used support-weighted F1 and cannot be relabelled as macro-F1.

Compared with GPT-5.1 alone, RetinAgent's paired macro-F1 differences were
0.3420 (95\% CI, 0.2723--0.4106) for glaucoma and 0.4432
(0.3823--0.5031) for DR. Corresponding Holm-adjusted McNemar $p$ values
were $1.44\times10^{-8}$ and $1.37\times10^{-15}$. The AMD difference
was 0.2237 (0.1590--0.2907; adjusted $p=9.98\times10^{-8}$), conditional
on the 210 valid agent responses. These comparisons assess the saved
configurations; they do not isolate the contribution of each agent component
or quantify language-model run-to-run variability.

% TABLE 1, page 8: replace with table1_recomputed.tex in the generated bundle.
% Supplementary intervals: table1_intervals.tex and paired_comparisons.tex.
% Do not retain bold best-value cells from the previous, differently defined table.

% DISCUSSION, page 13: replace the FairVision highest-performance assertion.
RetinAgent had higher macro-F1 point estimates than RETFound on the complete
glaucoma and DR subsets, but uncertainty in the paired differences prevents
a conclusion of superiority. Its advantage over GPT-5.1 alone was clearer
on these saved configurations. AMD results were conditional on valid
responses, and missing predictions and unresolved baseline-source versions
limit their interpretation. These analyses support further evaluation of
reliability-informed assistance, not a general claim that agentic arbitration
outperforms the strongest foundation model or reduces subgroup disparities.

% DISCUSSION, page 15: replace the conclusion's FairVision improvement clause.
Treating reliability as an explicit input to agentic reasoning yielded
promising point estimates in these retrospective comparisons, but an advantage
over the strongest foundation baseline was not established on the complete
FairVision subsets. Independent evaluation and resolution of incomplete
response cohorts are needed before stronger performance claims can be made.

% METHODS 4.2, page 16: replace the FairVision cohort/metric sentences only.
% Keep separate GDP and external descriptions for their own provenance review.
The locked FairVision cohorts contained 250 selected cases per task:
124 positive and 126 negative references for glaucoma, and 125 of each
class for AMD and DR. Saved selection files support sampling seed 2026
for AMD and DR. The seed and sampling procedure for the particular glaucoma
cohort could not be established from the saved records; the separate
seed-2026 glaucoma selection is not the same cohort.
Metrics were recomputed from valid binary reference/prediction pairs.
Forty AMD agent predictions and one glaucoma MIRAGE prediction were missing
or invalid. They were not imputed as negative, counted as escalations,
or replaced with predictions from another run. Valid-response counts are
reported per row, and incomplete rows do not estimate full-cohort performance.
One glaucoma reference was recovered from the dataset summary; its original
NPZ has not been independently verified.

Macro-F1 is the mean of the positive- and negative-class F1 scores with
fixed labels $\{0,1\}$ and zero for undefined class F1. It is distinct from
support-weighted F1 and positive-class F1. For FairVision, worst-group F1
is the minimum fixed-label macro-F1 across race, sex and age ($<50$, 50--69,
$\geq70$), excluding unknown demographic values, with no support cutoff.
Ethnicity is not included in this particular minimum. Single-class groups
are retained; a single-class group classified entirely correctly has
fixed-label macro-F1 of 0.5. This convention does not measure sensitivity
in a group without positive cases. Evaluation age bins do not establish
which age bins a historical reliability lookup used.

\paragraph{Statistical analysis.}
We used 10,000 class-stratified percentile bootstrap replicates (seed 20261005)
to estimate marginal 95\% confidence intervals for F1, balanced accuracy and
paired differences. Sensitivity and specificity intervals in the tables use
the exact Clopper--Pearson binomial method, avoiding degenerate percentile
intervals when all observed cases in a class were classified correctly.
The underlying bootstrap intervals are also retained in the exported data.
Cases were sampled with
replacement separately within positive and negative reference classes.
For comparisons, each sampled case retained both methods' predictions;
only common valid cases were included. The sampling unit was a case, not
a verified independent patient cluster. Intervals condition on the observed
class counts and fixed predictions, and do not include uncertainty from
language-model reruns, training, prompt development or missing responses.
Worst-group minima were recomputed across the fixed group set in every
replicate rather than tracking only the observed worst group. If a group
was absent in any replicate, no interval for that minimum was reported;
replicates were not discarded. Small subgroup support, ties between minima,
and boundary estimates can make bootstrap intervals unstable or degenerate.

Exact two-sided McNemar tests assessed paired error rates, not macro-F1.
Holm adjustment covered seven retrospective comparisons: RetinAgent versus
RETFound for glaucoma and DR; versus VisionFM OCT for AMD; versus GPT-5.1
for each task; and versus the manuscript-named URFound OCT for AMD as a
source-discrepancy sensitivity analysis. AMD comparisons used only valid
responses. In addition to marginal intervals, we exported Bonferroni
percentile intervals for the seven macro-F1 differences, using individual
confidence level $1-0.05/7$ (99.2857\%) for an approximate familywise 95\%
level. This correction does not cover every metric, subgroup, or earlier
model-selection exercise. Neither correction removes retrospective
comparator-selection effects or previous test-result inspection. Intervals
that include zero do not demonstrate either superiority or equivalence.

% Methods 4.4 placeholder, page 17: only the FairVision agent rule established
% by this audit can be stated here. Do not invent a universal LLM retry rule.
In the saved task-specific FairVision agent results, invalid or unparseable
diagnoses could be recorded as $-1$ and were excluded by the historical metric
helper. The revised analysis explicitly reports the resulting denominators
and missing predictions. This observation does not establish a uniform
retry or failure-handling policy across all language-model workflows.
